% chktex-file 1
% chktex-file 2
% chktex-file 3
% chktex-file 8
% chktex-file 11
% chktex-file 12
% chktex-file 24
% chktex-file 26
% chktex-file 36
% chktex-file 44

\chapter{电力电子器件}

前两章中，我们分别探讨了在理想开关与有一定损耗参数的开关控制之下，电路的计算方法，但是此类“开关”是一个抽象的模型，并不是电路中的实际元件。

曾经，小电压或小电流控制的电动开关常以机械继电器的形态出现，但在对开关频率要求较高时，其在高压下易拉弧并造成损耗，不利于快速控制，寿命也受限。因此不太可能依靠机械开关完成所讨论电路的快速通断。

在半导体工业的蓬勃发展之下，半导体器件的性能不断增强，在高频开关条件下几乎完全能够代替机械元件，因此电力电子行业中主要通过半导体器件充当开关。当然，半导体器件也具有各种特性，有的会带来损耗，故本章专门针对半导体器件的相关特性进行分析。

\section{从“开关”模型到具体器件}

以 Buck 为例，可以把一个单刀双掷（SPDT）开关分解为一对互斥的单刀单掷（SPST）开关 $\mathrm{A}$、$\mathrm{B}$。
\vspace{-0.3cm}

\InsertFigure{pic0301}

$\mathrm{A}$、$\mathrm{B}$ 除了互斥，还应满足各自的通态电流和断态电压要求：$\mathrm{A}$ 导通时承载 $i_{\mathrm{A}}(t)$，截止时承受 $v_{\mathrm{A}}(t)$；$\mathrm{B}$ 导通时承载 $i_{\mathrm{B}}(t)$，截止时承受 $v_{\mathrm{B}}(t)$。

那么，用哪些器件最适合呢？

\InsertFigure{pic0302}

部分器件自带体二极管或外接反并联二极管，绘图时可一并表示。

\subsection{器件的象限划分}

暂不考虑损耗，只考虑器件功能，即其允许的断态截止电压、通态通过电流方向，以及通断是否可控。

\InsertFigure{pic0303}

按照通态电流、断态电压划分象限，器件可分为单象限器件、电压双向器件、电流双向器件、四象限器件；象限定义具有相对性，笔记中以上图给出的参考方向为准。

\subsection{器件的控制能力}

按控制能力，分为全控、不控、半控。

不控器件的通断由外部电路条件决定；半控器件通常只能控制开通，而关断需由外部电路条件决定；全控器件可控制开通和关断。

\subsection{器件及其组合}

\subsubsection{二极管}

二极管为正向导通，反向阻断的不控型器件。由于通态电流与上文参考方向相同，而断态（截止）电压方向相反，故属于 II 象限器件。

\InsertFigure{pic0304}

\subsubsection{三极管、IGBT 与 MOSFET}

三极管、IGBT 为全控 I 象限器件。MOSFET 可阻断单一极性的电压，导通时沟道可双向通流，正向导通电流可受控关断；但由于体二极管（未画出）的存在，沟道关断后仍可反向通流，归为 I、IV 象限器件。

\InsertFigure{pic0305}

\subsubsection{二极管、三极管组合}

带反并联二极管的三极管如图 \ref{pic0306a} 所示，为 I、IV 象限；反向电流若经反并联二极管流动，不能由三极管的控制端直接关断。带串联二极管的三极管如图 \ref{pic0306b} 所示，为 I、II 象限。

\InsertFigure{pic0306}

\subsubsection{四象限器件组合}

四象限开关兼具双向通流和双向阻断能力，图 \ref{pic0307} 给出了三种实现方式。

\InsertFigure{pic0307}

图 \ref{pic0307a} 将两条由二极管与三极管串联的支路反并联，分别通过控制对应的三极管实现通断。两只三极管均关断时，组合可阻断两种极性的端口电压。

图 \ref{pic0307b} 将两只三极管共发射极连接，并分别配置反并联二极管。其正、反向电流分别由两只三极管控制，两只三极管均关断时，同样可实现双向阻断。

图 \ref{pic0307c} 利用二极管全桥，使外部电流无论沿哪个方向流动，通过三极管的电流方向均保持不变。只需控制一只三极管，即可控制两个方向的电流或阻断两个方向的电压。

\newpage{}

\section{二极管}

\subsection{器件特性与种类}

本节不展开二极管微观原理，只了解其端口特性。二极管的静态伏安特性如图 \ref{pic0308a} 所示，正向偏置时导通，反向偏置时阻断；反向电压超过击穿电压后进入击穿区。

\InsertFigure{pic0308}

普通 PN 结功率二极管属于双极型、少子器件，可用于较高电压和较大功率的整流场合。肖特基二极管属于单极型、多子器件，硅基器件通常用于较低电压场合，而 SiC 器件也可用于较高电压场合。

\subsection{器件开通与关断过程}

实际器件并非理想开关，下图为二极管开通与反向恢复波形。

\InsertFigure{pic0309}

普通 PN 结功率二极管开通时，少子积累与电导调制的建立需要时间，端口电压可能先出现过冲，再降至稳定的正向压降。关断时，电流先从正向电流降至零，再出现短暂的反向电流，随后恢复反向阻断能力。

\subsection{器件主要参数}

功率二极管的常用参数及其含义如下。

反向重复峰值电压 $V_{\mathrm{RRM}}$：在规定条件下，二极管允许重复承受的反向电压峰值上限。

额定正向平均电流 $I_{\mathrm{F(AV)}}$：在规定电流波形、温度和散热条件下，二极管允许通过的正向电流平均值上限。

非重复正向浪涌电流 $I_{\mathrm{FSM}}$：在规定脉冲波形和持续时间下，二极管能够短时承受的非重复正向电流峰值。

正向导通压降 $V_{\mathrm{F}}$：在规定正向电流和结温下，二极管导通时两端的电压，用于评估通态损耗。

反向漏电流 $I_{\mathrm{R}}$：在规定反向电压和结温下，二极管处于反向阻断状态时仍会流过的微小电流。

反向恢复时间 $t_{\mathrm{rr}}$：二极管由正向导通转为反向阻断时，从电流过零到反向恢复电流衰减至规定判据所需的时间。

反向恢复电荷 $Q_{\mathrm{rr}}$：反向恢复电流脉冲对时间积分的绝对值，表示反向恢复过程中流过的电荷量，与换流损耗有关。

\newpage{}

\section{MOSFET}

\subsection{器件特性}

本节以 N 沟道功率 MOSFET 为例，不展开其微观原理。其体二极管方向从源极指向漏极。当栅极关断受控沟道时，体二极管仍允许反向电流流过，某些情况下可起续流、保护作用；开通时，主要通过沟道实现双向导电。

\InsertFigure{pic0310}

\subsection{器件开通与关断过程}

作为开关使用时，MOSFET 的稳定工作状态为充分导通的线性区和截止区。MOSFET 开通与关断时，漏源电压和漏极电流会在一段时间内同时非零，形成开关损耗；稳定通态仍有通态电阻损耗。

\InsertFigure{pic0311}

忽略漏电与体二极管导通损耗时，MOSFET 及其栅极驱动的平均损耗分为三部分：
\InsertEquation{equ0301}
其中，$P_{\mathrm{loss}}$ 为总平均损耗，$P_{\mathrm{cond}}$、$P_{\mathrm{sw}}$、$P_{\mathrm{drv}}$ 分别为通态、开关和栅极驱动损耗。

通态损耗为
\InsertEquation{equ0302}
其中，$I_{\mathrm{D,rms}}$ 为漏极电流有效值，$R_{\mathrm{on}}$ 为通态电阻。

开关损耗为
\InsertEquation{equ0303}
其中，$f_{\mathrm{s}}$ 为开关频率，$E_{\mathrm{on}}$、$E_{\mathrm{off}}$ 分别为单次开通、关断的损耗能量。

栅极采用 $0\sim V_{\mathrm{GS}}$ 单极性驱动时，驱动损耗近似为
\InsertEquation{equ0304}
其中，$V_{\mathrm{GS}}$ 为栅源驱动电压（高电平），$Q_{\mathrm{g}}$ 为总栅极电荷。

\subsection{器件主要参数}

功率 MOSFET 的常用参数及其含义如下。

额定漏源电压 $V_{\mathrm{DSS}}$：栅源短接时，器件在规定条件下允许承受的最大漏源电压，用于确定关断耐压。

额定连续漏极电流 $I_{\mathrm{D}}$：在规定栅压、温度和散热条件下，器件允许连续通过的漏极电流上限。

最大脉冲漏极电流 $I_{\mathrm{DM}}$：在规定脉宽、占空比和温度条件下，器件允许短时通过的漏极电流峰值。

最大栅源电压 $V_{\mathrm{GS,\max}}$：栅源间允许施加的电压限值，正、负方向的额定范围以数据手册为准。

栅源阈值电压 $V_{\mathrm{GS(th)}}$：漏极电流达到规定微小测试电流时的栅源电压，表示开始形成导电沟道的条件，不能作为充分导通的驱动电压。

通态电阻 $R_{\mathrm{on}}$：数据手册常记为 $R_{\mathrm{DS(on)}}$，表示规定栅压、电流和结温下，器件导通时漏源间的等效电阻，用于计算通态损耗。

总栅极电荷 $Q_{\mathrm{g}}$：在规定漏压和漏极电流下，将栅极从关断状态驱动至给定栅压所需的总电荷，用于评估驱动需求与驱动损耗。

输入电容 $C_{\mathrm{iss}}$：$C_{\mathrm{iss}}=C_{\mathrm{gs}}+C_{\mathrm{gd}}$，反映输入端的电容充放电特性。

输出电容 $C_{\mathrm{oss}}$：$C_{\mathrm{oss}}=C_{\mathrm{ds}}+C_{\mathrm{gd}}$，反映输出端的电容充放电特性，与开关过程中的储能有关。

反向传输电容 $C_{\mathrm{rss}}$：$C_{\mathrm{rss}}=C_{\mathrm{gd}}$，表示栅漏间的反馈电容，与米勒效应有关。上述电容值均对应数据手册规定的测试条件。

常用的器件性能指标为
\InsertEquation{equ0305}
该指标反映通态电阻与栅极电荷的折中。在耐压等级、器件技术和测试条件相近时，数值较小通常更有利，但仍需结合输出电容及体二极管恢复特性比较。

\newpage{}

\exampletitle[exa0301]{计及开关过渡与二极管反向恢复的 Buck 变换器}

仍取 CCM、周期稳态和小纹波近似，仅计及开通过渡与二极管反向恢复，忽略通态压降、电感铜损及其他损耗，并将 $t_{\mathrm{rr}}$、$Q_{\mathrm{rr}}$ 视为固定参数。该简化模型用于说明开通过渡与反向恢复对变换比、输入电流和效率的影响。

\InsertFigure{pic0312}

\examplesubheading{器件分析}

图 \ref{pic0312} 中的 $\mathrm{Q}$ 开通时，$\mathrm{D}$ 先反向恢复，$i_{\mathrm{D}}<0$，故 $i_{\mathrm{Q}}=i_{\mathrm{L}}-i_{\mathrm{D}}>I_{\mathrm{L}}$；恢复后 $v_{\mathrm{d}}$ 升至 $V_{\mathrm{in}}$，如图 \ref{pic0313} 所示。

\InsertFigure{pic0313}

本例将开通过渡简化为持续 $t_{\mathrm{rr}}$ 的反向恢复阶段，其余开关延迟及电压跳变时间忽略。控制信号占空比为 $D_{\mathrm{c}}$，节点电压为高电平的有效占空比为 $D$，有
\InsertEquation{equ0306}

\examplesubheading{等效电路}

按 $v_{\mathrm{d}}$ 的高、低电平，画出图 \ref{pic0314} 所示的两种输出回路等效状态，持续时间分别为 $DT_{\mathrm{s}}$、$D'T_{\mathrm{s}}$。等效电路与图 \ref{pic0103} 相同，仍按节点电压判定状态。反向恢复期间虽已发出开通信号，输出回路仍按状态 2 处理；$\mathrm{Q}$ 的附加输入电流由图 \ref{pic0313} 单独计入。

\InsertFigure{pic0314}

共有 2 个等效状态与 2 个储能元器件，因此共 4 条方程。

\examplesubheading{状态方程}
\examplesubheading{小纹波近似}

接状态 1 时，
\InsertEquation{equ0307}

\examplesubheading{变换比与平均状态量}

由伏秒平衡，
\InsertEquation{equ0309}

输入平均电流由图 \ref{pic0313} 中 $i_{\mathrm{Q}}$ 的波形面积求得：
\InsertEquation{equ0315}

\examplesubheading{电感电流与电容电压纹波}

纹波波形同图 \ref{pic0104}、图 \ref{pic0106}，其中 $D$ 取有效占空比。沿用前章推导，
\InsertEquation{equ0316}

\examplesubheading{建立等效电路模型}

式 \eqref{equ0309} 和式 \eqref{equ0312} 分别给出输出端口的 KVL、KCL，式 \eqref{equ0315} 补充输入端口的 KCL。分别用受控源表示，得图 \ref{pic0315}\subref{pic0315a}；再将 $DV_{\mathrm{in}}$ 与 $DI_{\mathrm{L}}$ 合并为直流变压器，得图 \ref{pic0315}\subref{pic0315b}。 $Q_{\mathrm{rr}}/T_{\mathrm{s}}$ 和 $t_{\mathrm{rr}}I_{\mathrm{L}}/T_{\mathrm{s}}$ 为额外电流，保留原样。

\InsertFigure{pic0315}

\examplesubheading{效率与高效率条件}

在上述直流平均模型中，
\InsertEquation{equ0317}
将 $f_{\mathrm{s}}=1/T_{\mathrm{s}}$ 与式 \eqref{equ0314} 代入，得
\InsertEquation{equ0320}
要使效率接近 1，应满足
\InsertEquation{equ0321}

\InsertFigure{pic0316}

图 \ref{pic0316} 仅表示 $t_{\mathrm{rr}}$、$Q_{\mathrm{rr}}$ 固定时的模型趋势，当 $t_{\mathrm{rr}}=Q_{\mathrm{rr}}=0$ 时，退化为前章的理想 Buck 关系。

\newpage{}

\section{BJT 与 IGBT}

\subsection{器件特性}

BJT 即三极管，是电流控制型器件，基极、集电极、发射极分别记为 $\mathrm B$、$\mathrm C$、$\mathrm E$ ，导通期间需要持续提供基极电流。

\InsertFigure{pic0317}

IGBT 是绝缘栅双极型晶体管，栅极、集电极、发射极分别记为 $\mathrm G$、$\mathrm C$、$\mathrm E$，通过栅极、发射极间电压控制开通与关断。

\InsertFigure{pic0318}

\subsection{器件开通与关断过程}

BJT 的驱动电路与开关波形如图 \ref{pic0319} 所示。

\InsertFigure{pic0319}

BJT 开通时，基极电流使集电极电流上升、集电极电压下降；关断时需移除存储电荷，反向基极电流可加快这一过程。开关过程中的电压、电流重叠会产生损耗。

IGBT 关断时，沟道先关断，但存储少子不能立即消失，因此产生电流拖尾，如图 \ref{pic0320} 所示。

\InsertFigure{pic0320}

图中，阶段 (1) 为集电极电压上升；阶段 (2) 为集电极电流快速下降，电流逐渐转移至二极管；阶段 (3) 为电流拖尾。拖尾期间 IGBT 已承受较高电压，仍有功率损耗。

将导通压降近似为常数时，BJT 与 IGBT 的导通损耗可写为
\InsertEquation{equ0322}
其中 $I_{\mathrm{C,av}}$ 为一个开关周期内的集电极电流平均值。

\subsection{器件主要参数}

BJT 与 IGBT 的常用参数及其含义如下。

集电极、发射极额定电压：BJT 常用 $V_{\mathrm{CEO}}$，表示基极开路时允许承受的集电极、发射极电压上限；IGBT 常用 $V_{\mathrm{CES}}$，表示栅极、发射极短接时相应的电压上限。

额定集电极电流 $I_{\mathrm C}$：在规定温度、驱动和散热条件下，器件允许连续通过的集电极电流上限。

脉冲集电极电流 $I_{\mathrm{CM}}$：在规定脉宽、占空比和结温条件下，器件允许短时通过的集电极电流峰值。

集电极、发射极饱和压降 $V_{\mathrm{CE(sat)}}$：在规定集电极电流、驱动和结温下，器件导通时的压降，用于评估通态损耗。

直流电流增益 $h_{\mathrm{FE}}$：BJT 在规定放大区工作条件下，集电极电流与基极电流的比值。饱和驱动时采用的强迫电流增益 $I_{\mathrm C}/I_{\mathrm B}$ 应与其区分。

栅极、发射极极限电压 $V_{\mathrm{GES}}$：IGBT 栅极、发射极之间允许施加的电压上限，正、负极限应分别按数据手册取值。

栅极阈值电压 $V_{\mathrm{GE(th)}}$：IGBT 在规定小电流条件下开始导通时的栅极、发射极电压，不代表充分导通所需的驱动电压。

栅极电荷 $Q_{\mathrm g}$：IGBT 在规定工作条件下完成栅极充电所需的总电荷量，与驱动能力及驱动损耗有关。

开关时间：BJT 主要包括延迟时间 $t_{\mathrm d}$、上升时间 $t_{\mathrm r}$、存储时间 $t_{\mathrm{stg}}$ 和下降时间 $t_{\mathrm f}$；IGBT 常列出开通、关断延迟时间及电流上升、下降时间。具体起止判据应按数据手册取值，$t_{\mathrm f}$ 不能直接代替完整关断时间。

单次开关能量 $E_{\mathrm{on}}$、$E_{\mathrm{off}}$：在规定电压、电流、温度和驱动条件下，开通、关断过程的功率积分，用于计算开关损耗；IGBT 的 $E_{\mathrm{off}}$ 包括规定积分范围内的拖尾损耗。

\newpage{}

\section{SCR}

\subsection{器件特性}

普通单向晶闸管（SCR）是半控型器件，阳极、阴极、门极分别记为 $\mathrm A$、$\mathrm K$、$\mathrm G$。门极可控制开通，导通并锁定后撤去门极信号不能使其关断。

\InsertFigure{pic0321}

图 \ref{pic0321b} 中，$\mathrm Q_1$ 为 PNP 管，$\mathrm Q_2$ 为 NPN 管，两管的集电极分别接至对方基极，形成正反馈。门极触发后，两管相互增强，使器件由阻断转入导通。

SCR 的静态伏安特性如图 \ref{pic0322} 所示。正向偏置时可由门极触发导通；反向偏置时阻断，超过反向击穿电压后进入击穿区。$V_{\mathrm{BO}}$ 为门极开路时的正向转折电压，门极触发可使器件在较低的正向电压下导通。

\InsertFigure{pic0322}

\subsection{器件开通与关断过程}

开通时，在正向偏置的 SCR 门极施加正向触发电流。阳极电流达到锁定电流 $I_{\mathrm L}$ 后，即使撤去门极信号，器件仍能保持导通。

关断时，需要使阳极电流降至维持电流 $I_{\mathrm H}$ 以下。电流过零后，还需在规定反向偏置条件下留出不小于关断时间 $t_{\mathrm q}$ 的恢复时间，才能重新承受正向电压。交流电路可利用电流自然过零换流，直流电路则需由外部电路创造关断条件。

\newpage{}

\subsection{器件主要参数}

SCR 的常用参数及其含义如下。

断态重复峰值电压 $V_{\mathrm{DRM}}$ 与反向重复峰值电压 $V_{\mathrm{RRM}}$：在规定条件下，器件允许重复承受的正向阻断电压和反向电压峰值上限。

额定通态平均电流 $I_{\mathrm{T(AV)}}$：在规定波形、温度和散热条件下，器件允许通过的通态电流平均值上限；数据手册也常给出通态电流有效值上限 $I_{\mathrm{T(RMS)}}$。

非重复通态浪涌电流 $I_{\mathrm{TSM}}$：在规定脉冲波形和持续时间下，器件能够短时承受的非重复通态电流峰值。

通态峰值电压 $V_{\mathrm{TM}}$：在规定通态峰值电流和结温下，器件导通时两端的电压，用于评估通态损耗。

门极触发电流 $I_{\mathrm{GT}}$ 与门极触发电压 $V_{\mathrm{GT}}$：在规定条件下，使器件由正向阻断转为导通所需的门极电流及相应的门极、阴极电压。

锁定电流 $I_{\mathrm L}$：开通后撤去门极触发脉冲时，能使器件保持导通的最小阳极电流。

维持电流 $I_{\mathrm H}$：门极无驱动时，器件维持已建立的导通状态所需的最小阳极电流，通常小于锁定电流。

关断时间 $t_{\mathrm q}$：在规定换流条件下，阳极电流过零后，器件恢复正向阻断能力所需的最小时间。

临界电压上升率 $(\mathrm{d}v/\mathrm{d}t)_{\mathrm{cr}}$ 与允许电流上升率 $(\mathrm{d}i/\mathrm{d}t)_{\mathrm{cr}}$：分别限制断态电压上升过快引起的误开通，以及开通过程中电流上升过快引起的局部过热。

\newpage{}

\section{宽禁带器件}

宽禁带功率器件主要采用 SiC、GaN 等材料。SiC 器件适用于高压、较大功率场合，在部分场合正在逐渐替代 IGBT 以提高电路工作频率。GaN 器件常为 HEMT 结构，可在部分高频电源中替代硅 MOSFET。

\section{器件应用}

\subsection{各类器件的应用范围}

\InsertFigure{pic0323}

\InsertFigure{pic0324}

\clearpage

\section{第三章符号汇总}

本节汇总本章主要符号。电路量中，大写表示直流分量或平均值，小写表示瞬时量；器件参数按本章定义。$x$ 和 $y$ 表示任意字符串。未列举的符号不是主要符号。

\InsertTable{tab0301}

\clearpage

\InsertTable{tab0302}

